\chapter{Linear Programming in Geometry}
\label{ch:linear_programming}

Linear programming provides a common language for geometric optimization.
Many geometric problems ask for a point satisfying linear inequalities while
maximizing or minimizing a linear objective. In two dimensions, the feasible
region is an intersection of half-planes and therefore a convex polygon.

\section{Geometric Linear Programs}
\label{sec:geometric_linear_programs}

A linear program has the form
\[
  \text{maximize } c^{\mathsf T}x
  \qquad\text{subject to}\qquad Ax\leq b.
\]
Each row of $Ax\leq b$ defines a half-plane. The feasible set is convex, and
when it is non-empty and bounded, an optimum occurs at a vertex of the
feasible polytope. In two dimensions this reduces optimization to a convex
polygon; in three dimensions it reduces to a convex polyhedron.

Applications include smallest enclosing regions, collision separation,
supporting-line queries, facility placement, and motion-planning constraints.

\section{Half-Plane Intersection}
\label{sec:halfplane_intersection}

For an oriented line through points $p$ and $q$, the left half-plane is
\[
  H(p,q)=\{x\in\mathbb{R}^2:\operatorname{orient2d}(p,q,x)\geq 0\}.
\]
The intersection of a finite collection of half-planes is a convex polygon,
possibly empty or unbounded. A bounding box can be added when a finite output
polygon is required.

One robust approach sorts the boundary lines by direction and maintains a
deque of candidate lines. Whenever the intersection of the last two lines is
outside the next half-plane, the last line is removed. A symmetric check is
performed at the front of the deque.

\begin{algorithm}[htbp]
\caption{Half-Plane Intersection}
\label{alg:halfplane_intersection}
\begin{algorithmic}[1]
\Function{HalfPlaneIntersection}{$H$}
  \State Sort half-planes by boundary direction
  \State $D\gets$ empty deque
  \For{$h$ in $H$}
    \While{$|D|\geq 2$ and intersection of the last two lines is outside $h$}
      \State Remove the last half-plane from $D$
    \EndWhile
    \While{$|D|\geq 2$ and intersection of the first two lines is outside $h$}
      \State Remove the first half-plane from $D$
    \EndWhile
    \State Append $h$ to $D$
  \EndFor
  \While{$|D|\geq 3$ and the last intersection is outside the first half-plane}
    \State Remove the last half-plane from $D$
  \EndWhile
  \While{$|D|\geq 3$ and the first intersection is outside the last half-plane}
    \State Remove the first half-plane from $D$
  \EndWhile
  \State \Return polygon formed by consecutive boundary intersections
\EndFunction
\end{algorithmic}
\end{algorithm}

After sorting, the deque algorithm takes $O(n)$ time and $O(n)$ space. Parallel
or coincident boundary lines require a consistent rule: retain the more
restrictive half-plane when their feasible sides agree, and report an empty
intersection when they conflict.

\section{Low-Dimensional Linear Programming}
\label{sec:low_dimensional_lp}

In fixed dimension, randomized incremental algorithms solve linear programs
efficiently by processing constraints in random order. When a new constraint
violates the current optimum, the optimum must lie on that constraint's
boundary, reducing the problem dimension by one. This gives expected linear
time for two-dimensional linear programming and expected linear time for
fixed-dimensional LP under the usual boundedness and constant-dimension
assumptions.

The same boundary-reduction idea appears in smallest-enclosing-circle
algorithms, support-point queries, and randomized convex-hull construction.
It is therefore a useful bridge between predicates, convexity, optimization,
and geometric data structures.

\section{Duality and Applications}
\label{sec:lp_duality_applications}

Point-line duality maps a point $(a,b)$ to a line such as $y=ax-b$. Under a
consistent duality transform, incidence and above/below relationships are
preserved or reversed in a controlled way. This converts convex-hull and
half-plane questions into upper- or lower-envelope questions.

Linear constraints also describe collision-free configuration-space regions,
visibility restrictions, separating planes, and feasible placement regions.
Numerical implementations should use the robust orientation predicates from
Chapter~\ref{ch:geometric_predicates} when testing boundary intersections.

